% Unit: Computer Science A --- Spring Core/Boot, Web/MVC, Data JPA, Actuator, Lombok (\input after \texttt{cs-a/05-practice.tex} in \texttt{main.tex}).

\runningheader{Computer Science A}{Unit 6: Spring ecosystem}{Page \thepage\ of \numpages}

\begingradingrange{csau6}

\uplevel{\par\textbf{Spring core: IoC and beans}\par}

\question[4]
What is Inversion of Control? What are its advantages?
\makeemptybox{1.45in}
\begin{solution}
  Framework calls your code and supplies dependencies instead of your code \emph{new}-ing collaborators. Advantages: looser coupling, easier testing (inject mocks), centralized wiring, and clearer separation of concerns.
\end{solution}

\question[4]
What is Dependency Injection?
\makeemptybox{1.2in}
\begin{solution}
  Objects receive collaborators from outside (constructor/setter/field injection) rather than creating them internally; Spring's IoC container resolves and injects beans.
\end{solution}

\question[4]
What is the difference between XML-based and annotation-based configuration?
\makeemptybox{1.35in}
\begin{solution}
  XML externalizes bean definitions in files; annotations (\texttt{@Component}, \texttt{@Configuration}, \texttt{@Bean}) colocate metadata with code. Modern Spring favors Java config and annotations with optional XML for legacy or cross-cutting wiring.
\end{solution}

\question[4]
What is a Bean?
\makeemptybox{1.05in}
\begin{solution}
  In Spring, a managed object instantiated, assembled, and lifecycle-managed by the IoC container, identified by name/type and wired per configuration.
\end{solution}

\question[4]
What is the Spring Framework? What are its benefits?
\makeemptybox{1.35in}
\begin{solution}
  Mature Java ecosystem for enterprise apps: IoC/DI, AOP, data access, web stack, integration. Benefits include portability, batteries-included modules, large community, and consistent programming model.
\end{solution}

\question[4]
What does the Spring Core module provide?
\makeemptybox{1.2in}
\begin{solution}
  IoC container, bean lifecycle, \texttt{ApplicationContext}, resource abstraction, SpEL, and foundational utilities other modules build on.
\end{solution}

\question[4]
How could you create an IoC Container?
\makeemptybox{1.2in}
\begin{solution}
  Use an \texttt{ApplicationContext} implementation: Java-config or XML context classes, or call \texttt{SpringApplication.run} in Boot to build the context from your configuration and component scan.
\end{solution}

\newpage

\uplevel{\par\textbf{Spring Boot and configuration}\par}

\question[4]
What is Spring Boot?
\makeemptybox{1.25in}
\begin{solution}
  Opinionated layer on Spring: dependency starters, auto-configuration, embedded HTTP servers, and sensible defaults so apps start with little boilerplate.
\end{solution}

\question[4]
Tell me about the \texttt{@SpringBootApplication} annotation.
\makeemptybox{1.35in}
\begin{solution}
  Composes \texttt{@Configuration}, \texttt{@EnableAutoConfiguration}, and \texttt{@ComponentScan}: marks the main class, enables Boot auto-config, and scans for Spring components in the package hierarchy.
\end{solution}

\question[4]
What is Spring Initializr?
\makeemptybox{1.1in}
\begin{solution}
  Web service/UI (start.spring.io) that generates a Boot project skeleton with chosen dependencies, build tool, and Java version.
\end{solution}

\question[4]
What is the use of a \texttt{.properties} or \texttt{.yaml} file?
\makeemptybox{1.35in}
\begin{solution}
  Externalize configuration (ports, URLs, credentials via env, feature flags) without recompiling; Boot binds prefixed keys to beans and supports profiles (\texttt{application-dev.yml}).
\end{solution}

\newpage

\uplevel{\par\textbf{Spring Web and MVC}\par}

\question[4]
Describe MVC architecture.
\makeemptybox{1.35in}
\begin{solution}
  Model holds data/state, View renders output (or API representation), Controller handles input, delegates to services, and chooses view/response; Spring MVC maps HTTP to controllers with clear separation.
\end{solution}

\question[4]
Describe commonly used Spring Web annotations.
\makeemptybox{1.55in}
\begin{solution}
  Map URLs to methods with mapping annotations; bind inputs with parameter annotations; mark REST controllers; set status and bodies with response helpers.
\end{solution}

\question[4]
What does the Spring Web module provide?
\makeemptybox{1.2in}
\begin{solution}
  DispatcherServlet, handler mapping, view resolution (for MVC), REST support, HTTP message conversion, filters, CORS, and integration with servlet containers or reactive stacks.
\end{solution}

\question[4]
What are some annotations that can be used to get information from the request?
\makeemptybox{1.35in}
\begin{solution}
  \texttt{@RequestParam}, \texttt{@PathVariable}, \texttt{@RequestHeader}, \texttt{@CookieValue}, \texttt{@RequestBody}, and sometimes \texttt{HttpServletRequest} parameters for low-level access.
\end{solution}

\question[4]
What can we use to customize the response sent back to the client?
\makeemptybox{1.25in}
\begin{solution}
  Return types (\texttt{ResponseEntity}), \texttt{@ResponseBody}/\texttt{@RestController}, HTTP status annotations, message converters, headers, and exception handlers that map errors to bodies/status.
\end{solution}

\question[4]
How does the \texttt{@ExceptionHandler} annotation work?
\makeemptybox{1.3in}
\begin{solution}
  Methods in a \texttt{@Controller}/\texttt{@RestController} (or \texttt{@ControllerAdvice}) handle specific exception types, build an error response, and avoid generic 500 pages for API clients.
\end{solution}

\newpage

\uplevel{\par\textbf{Spring Data JPA}\par}

\question[4]
What configuration is needed in the \texttt{.properties} or \texttt{.yaml} file to use the Spring Data module?
\makeemptybox{1.35in}
\begin{solution}
  Data source keys (URL, user, password, driver) and JPA/Hibernate settings (DDL mode, dialect). Boot builds a \texttt{DataSource} and \texttt{EntityManagerFactory} from them.
\end{solution}

\question[4]
What is the JPA API? What does it provide?
\makeemptybox{1.3in}
\begin{solution}
  Java Persistence API: standard ORM mapping (\texttt{@Entity}, relationships), \texttt{EntityManager} for CRUD/query, and portability across providers (Hibernate, EclipseLink, etc.).
\end{solution}

\question[4]
What is an ORM?
\makeemptybox{1.05in}
\begin{solution}
  Object-Relational Mapping: maps classes to tables and object graphs to SQL, reducing hand-written JDBC for common persistence patterns.
\end{solution}

\question[4]
What is Spring Data JPA?
\makeemptybox{1.2in}
\begin{solution}
  Spring Data module that reduces repository boilerplate: derive queries from method names, \texttt{@Query}, paging/sorting, and integrates \texttt{EntityManager} with Spring transactions.
\end{solution}

\question[4]
How do you create a repository with Spring Data JPA?
\makeemptybox{1.2in}
\begin{solution}
  Declare an interface extending \texttt{JpaRepository<Entity, Id>} (or related) in a scanned package; Spring generates the implementation at runtime---no manual DAO class.
\end{solution}

\question[4]
What methods does the \texttt{JpaRepository} interface provide?
\makeemptybox{1.35in}
\begin{solution}
  Inherits CRUD and paging/sorting (\texttt{save}, \texttt{findById}, \texttt{deleteAll}, etc.) and adds JPA-oriented helpers like flush and batched persistence APIs (exact extras vary by Spring Data version).
\end{solution}

\question[4]
What are custom queries and give me an example of how you would write a custom method.
\makeemptybox{1.45in}
\begin{solution}
  JPQL or native SQL declared with \texttt{@Query} on a repository method, named queries, or \texttt{Specification}/Criteria when you cannot express the rule with derived method names alone.
\end{solution}

\question[4]
What is the \texttt{@Transactional} annotation? What are its benefits?
\makeemptybox{1.35in}
\begin{solution}
  Declares a transactional boundary (often on services): rollback on errors, configurable propagation and isolation, delegated to Spring's transaction manager instead of hand-written JDBC \texttt{commit}/\texttt{rollback}.
\end{solution}

\newpage

\uplevel{\par\textbf{Lombok and Actuator}\par}

\question[4]
What is Lombok and its benefits?
\makeemptybox{1.25in}
\begin{solution}
  Annotation processor that generates getters, constructors, \texttt{equals}/\texttt{hashCode}, builders, etc., cutting POJO noise; teams should enable IDE support for accurate navigation.
\end{solution}

\question[4]
What is the Spring Actuator module?
\makeemptybox{1.2in}
\begin{solution}
  Production endpoints for health, metrics, environment, logging levels, and integration with monitoring systems; complements Boot's operational defaults.
\end{solution}

\question[4]
What are actuator endpoints?
\makeemptybox{1.1in}
\begin{solution}
  HTTP or JMX operations under \texttt{/actuator} (by default) such as health and metrics; each has an id, can be enabled or exposed separately, and should be secured in production.
\end{solution}

\question[4]
What does enabling an endpoint mean?
\makeemptybox{1.05in}
\begin{solution}
  The endpoint exists in the application and can run its logic if invoked---controlled via \texttt{management.endpoint.<id>.enabled} properties.
\end{solution}

\question[4]
What does exposing an endpoint mean?
\makeemptybox{1.05in}
\begin{solution}
  Making an enabled endpoint reachable over a technology (commonly HTTP or JMX) via \texttt{management.endpoints.web.exposure.include} or related settings.
\end{solution}

\question[4]
What does it mean when an endpoint is ``available''?
\makeemptybox{1.1in}
\begin{solution}
  Boot considers it both enabled and exposed on the chosen interface so clients can call it subject to security rules.
\end{solution}

\question[4]
What are some endpoints that actuator provides and their use?
\makeemptybox{1.45in}
\begin{solution}
  Examples: \texttt{health} for load balancers, \texttt{info} for build metadata, \texttt{metrics} for observability, \texttt{env}/\texttt{configprops} for debugging configuration (lock down in production).
\end{solution}

\question[4]
How can you create a custom actuator endpoint?
\makeemptybox{1.25in}
\begin{solution}
  Implement \texttt{Endpoint} or \texttt{@Endpoint} / \texttt{@ReadOperation} (Boot 2+) beans, register with a unique id, secure them, and optionally expose over HTTP via management configuration.
\end{solution}

\endgradingrange{csau6}
